\section{Language models from the ground up}
\label{sec:lm}

\begin{sessioncard}
\textbf{Session 1.} \emph{Goal:} understand what a language model takes as input, what it produces as output,
what its weights are and how they are found, and which question interpretability tries to answer.
\emph{Experiments:} how GPT-2 small splits text into tokens (\cref{exp:tokens}), which tokens it predicts
(\cref{exp:predictions}), and a task that serves as the running example in all later sessions (\cref{exp:ioi}).
\emph{Prerequisites:} vectors, matrices, probabilities.
\end{sessioncard}

\subsection{What a language model does}

A language model is a function. Its input is a piece of text, called the \emph{prompt}. Its output is a
probability distribution over the next piece of text. A language model can answer questions and write programs,
but it does all of this by applying this one function many times. The model predicts the next piece of text. One
piece is drawn at random from this prediction and appended to the input. Then the model is applied again to the
longer text, and so on.

The pieces of text are called \emph{tokens}. A token is a frequent word or a part of a word, usually together
with the space in front of it. A fixed procedure, called the \emph{tokenizer}, splits every text into tokens. It
then replaces each token by its number in a fixed list of tokens, called the \emph{vocabulary}.

\begin{experiment}[Tokens]\label{exp:tokens} \script{exp\_basics.py}
The vocabulary of GPT-2 has $50{,}257$ tokens. Its tokenizer splits the sentence \tok{Interpretability is fun.}
into the six tokens
\begin{center}
\fbox{\tok{Inter}}\,\fbox{\tok{pret}}\,\fbox{\tok{ability}}\,\fbox{\tok{\textvisiblespace is}}\,\fbox{\tok{\textvisiblespace fun}}\,\fbox{\tok{.}}
\end{center}
where \tok{\textvisiblespace} marks a space. The rare word \tok{Interpretability} is split into three parts,
while the common word \tok{\textvisiblespace is} is a single token. The software also puts a special start token
in front of every text. We call this token BOS, which stands for ``beginning of sequence''.
\end{experiment}

At every position of the input, the model returns one probability for each of the $50{,}257$ tokens. These
probabilities are the model's prediction for the token at the next position. We will almost always look at the
last position, where the model predicts the token that follows the whole input.

\begin{experiment}[What a small model predicts]\label{exp:predictions} \script{exp\_basics.py}
Throughout these notes we study the model GPT-2 small \citep{radford2019language}. It was published in 2019, and
it runs on a laptop. \Cref{tab:predictions} shows its five most likely next tokens for three inputs. The model
continues the counting sequence correctly and with high probability. It completes the sentence about two people
with the correct name. For the question of general knowledge, it gives London and Paris the same small
probability. The last row teaches a lesson that is important for everything that follows. We can only ask
\emph{how} a model performs a task if the model actually performs it. Every study of a mechanism should therefore
start by checking that the model shows the behaviour in question.
\end{experiment}

\begin{table}[t]
\centering\small
\begin{tabular}{@{}>{\raggedright\arraybackslash}p{0.4\textwidth}>{\raggedright\arraybackslash}p{0.54\textwidth}@{}}
\toprule
Input & Most likely next tokens, with probability \\
\midrule
\tok{One, two, three, four,} & \tok{five} $0.85$, \tok{ten} $0.02$, \tok{Five} $0.01$, \tok{six} $0.01$, \tok{four} $0.01$ \\
\tok{When Mary and John went to the store, John gave a drink to} & \tok{Mary} $0.68$, \tok{them} $0.13$, \tok{the} $0.05$, \tok{John} $0.02$, \tok{her} $0.01$ \\
\tok{The Eiffel Tower is located in the city of} & \tok{London} $0.07$, \tok{Paris} $0.07$, \tok{Amsterdam} $0.04$, \tok{Berlin} $0.03$, \tok{New} $0.03$ \\
\bottomrule
\end{tabular}
\caption{Predictions of GPT-2 small for the token that follows three inputs (\cref{exp:predictions}).}
\label{tab:predictions}
\end{table}

\subsection{The building blocks of a neural network}

The function that turns tokens into probabilities is a \emph{neural network}. Three notions are enough to
describe a neural network: vectors, layers and weights.

\paragraph{Vectors.} All intermediate results that the network computes are vectors, that is, lists of real
numbers. For two vectors $x,w\in\R^d$, the inner product $w^\top x$ is the sum of the products of corresponding
entries. It measures how much $x$ points in the direction of $w$. A matrix $W_{\mathrm{in}}$ whose rows are
vectors of $\R^d$ turns a vector $x\in\R^d$ into the vector $W_{\mathrm{in}}x$. The entries of
$W_{\mathrm{in}}x$ are the inner products of $x$ with the rows of $W_{\mathrm{in}}$.

\paragraph{Layers.} A \emph{layer} is a map of the form
\begin{equation}\label{eq:layer}
x \;\longmapsto\; \relu(W_{\mathrm{in}}x+b_{\mathrm{in}}),
\end{equation}
with a matrix $W_{\mathrm{in}}$, a vector $b_{\mathrm{in}}$ called the \emph{bias}, and the function
$\relu(s)=\max(0,s)$, which is applied to every entry. Each entry of the result is called a \emph{neuron}, and its
value on a given input is called its \emph{activation}. So a neuron computes the inner product of its input with
one row of $W_{\mathrm{in}}$ and adds a number. It outputs the result if the result is positive, and zero
otherwise. The function $\relu$ is essential. Without it, several layers applied one after another would only
compute a function of the form $x\mapsto Ax+c$ with a matrix $A$ and a vector $c$. Such a function is far too
simple to predict text.

\paragraph{Weights.} A neural network applies many such layers one after another. \Cref{sec:transformer} describes
the variations of this construction that language models use. The entries of all matrices and bias vectors are
called the \emph{weights} or \emph{parameters} of the network. GPT-2 small has about $124$ million weights.
Current large models are built in the same way and have between $100$ and $1000$ times more weights.

\subsection{From scores to probabilities}

The last step of the network produces one real number for every token of the vocabulary. These numbers are
called \emph{logits}. We write $\logit(t)$ for the logit of token $t$ at the position that we consider. The
probabilities are computed from the logits by the \emph{softmax} function,
\begin{equation}\label{eq:softmax}
\Pr(\text{next token is }t) \;=\; \frac{\exp(\logit(t))}{\sum_{t'}\exp(\logit(t'))},
\end{equation}
where the sum runs over the whole vocabulary. For example, three tokens with the logits $2$, $1$ and $0$ receive
the probabilities $0.67$, $0.24$ and $0.09$.

We will use two properties of the softmax function many times. First, adding the same number to all logits does
not change the probabilities. Second, the ratio of the probabilities of two tokens $t$ and $t'$ equals
$\exp(\logit(t)-\logit(t'))$. Because of the first property, the value of a single logit has no meaning by itself.
Because of the second property, the difference of two logits has a clear meaning: it is the logarithm of how many
times more likely one token is than the other. \Cref{ex:softmax} asks for the proofs.

\subsection{How the weights are found}

Nobody chooses the weights of a language model by hand. The weights are found by \emph{training} on a large
collection of texts. Take a text from this collection and a position in it. The \emph{loss} at this position is
$-\log$ of the probability that the model assigns to the token that actually comes next. The loss is zero if the
model gives this token the probability one, and the loss is large if the model gives this token a small
probability. Training starts from random weights and repeats the following step many times. It computes the
average loss on a batch of texts, it computes the gradient of this average loss with respect to all weights, and
it moves the weights a small step in the direction opposite to the gradient.

Computing the output of the model for an input is called a \emph{forward pass}. Computing the gradient is called
a \emph{backward pass}. A backward pass costs only a few times as much as a forward pass, however many weights the
model has. After training, the weights are fixed, and the model is a deterministic function from tokens to
logits.

This is where the problem studied in these notes comes from. The training procedure is simple, and we understand
it fully. But the result of training is a list of $124$ million numbers. These numbers implement some algorithm for
predicting text, and no person designed this algorithm. The situation is similar to having the compiled binary
file of a computer program without its source code. We can see every instruction, but we cannot see the logic of
the program.

\subsection{The question of interpretability}

\begin{definition}[Interpretability]\label{def:interp}
An \emph{interpretation} of a model is a description of its internal computation that is (i) understandable by
a human and (ii) \emph{faithful}. Faithful means that the description correctly predicts what the model computes,
including what it computes after we change some of its internal values.
\emph{Mechanistic interpretability} looks for interpretations at the level of algorithms. It asks which
quantities the model represents, where in the model it represents them, and by which operations it combines them.
\end{definition}

The two requirements are in tension with each other. The full list of weights is a perfectly faithful
description, but no human can understand it. A one-sentence story is easy to understand, but it may be false.

\paragraph{Why not ask the model?} When we ask a modern model to explain one of its answers, it writes an
explanation. But this explanation is just more text generated by the same model. Nothing forces the explanation to
describe the computation that actually produced the answer. \citet{turpin2023language} tested this with
multiple-choice questions. In the prompt, they arranged the example questions so that the correct answer was
always the first option. The models then often chose the first option for a new question as well. They also wrote
a step-by-step justification of their choice, and this justification never mentioned the pattern in the examples
that had caused the choice. Such an explanation is \emph{plausible}, but it is not faithful
\citep{jacovi2020towards}. To obtain a faithful explanation, we have to look inside the model.

\paragraph{Three kinds of evidence.} The methods of mechanistic interpretability can be sorted by the kind of
evidence that they provide. It is important to keep these kinds of evidence apart, because they support claims of
different strength.
\begin{itemize}
\item \emph{Observational} methods record the activations of the model on some inputs and look for patterns in
them. These methods are cheap, but they can only show correlations. Sessions 2 to 5 contain several observational
methods.
\item \emph{Interventional} methods change an activation and measure how the output changes. These methods show
causal effects, but only for the particular change that was made. They first appear in session 3, and session 6
is devoted to them.
\item \emph{Mechanistic} explanations derive the behaviour of the model from its weights. Such an explanation also
covers inputs that were never tested. Session 8 presents a network for which this has been achieved.
\end{itemize}
A claim about a model should never be stronger than the evidence that supports it.

\subsection{The running example}
\label{sec:running}

A method is easier to understand when we see it applied to a concrete example. We therefore apply the methods of
these notes to one small task whenever this is possible. This task is called \emph{indirect object
identification} (IOI), and it was introduced by \citet{wang2023interpretability}. Consider again the second input
of \cref{tab:predictions}:
\begin{center}
\tok{When Mary and John went to the store, John gave a drink to}
\end{center}
A human reader continues with \tok{Mary}. To do so, the reader has to notice three things. Two names were
introduced, the name \tok{John} appears a second time, and the other name, \tok{Mary}, is the person who receives
the drink. The name that should be predicted is called the \emph{indirect object} (IO), and the repeated name is
called the \emph{subject} (S).

\begin{experiment}[GPT-2 small solves the task]\label{exp:ioi} \script{exp\_ioi.py}
We use 12 prompts of the same form as the sentence above. They are built from six pairs of names, and each pair
is used in both orders. On all 12 prompts, the most likely next token is the correct name. The average probability
of the correct name is $0.48$, and the average probability of the repeated name is $0.014$. We measure how well the
model performs on a prompt by the \emph{logit difference} $\logit(\mathrm{IO})-\logit(\mathrm{S})$ at the last
position. It is $3.70$ on average, and it lies between $2.06$ and $5.53$ on the 12 prompts. By the second property
of the softmax function, a logit difference of $3.70$ means that the correct name is $e^{3.70}\approx40$ times as
likely as the repeated name.
\end{experiment}

The model was trained only to predict the next token, and nobody taught it this task. How does it find the
correct name? The following sessions answer this question step by step, and session 7 puts the pieces together.

\begin{caution}
All numbers and data figures about GPT-2 small in these notes were produced by the scripts in the folder
\texttt{code/}. They illustrate the methods on small sets of prompts. Where they repeat an experiment from the
literature, they agree with the published results qualitatively. They are not meant as new research findings.
\end{caution}

\begin{takeaways}
A language model maps a sequence of tokens to logits, and the softmax function turns the logits into
probabilities for the next token. Only differences of logits have a meaning. The weights are found by gradient
descent on a prediction loss, so even the people who trained the model do not know the algorithm that the weights
implement. An interpretation has to be understandable and faithful. Evidence about the inside of a model can be
observational, interventional or mechanistic.
\end{takeaways}

\subsection*{Questions for the seminar}

\begin{exercise}[Softmax sees only differences]\label{ex:softmax}
(a) Show that adding the same number to all logits does not change the probabilities \eqref{eq:softmax}.
(b) Show that the ratio of the probabilities of two tokens $t$ and $t'$ equals $\exp(\logit(t)-\logit(t'))$.
(c) Let $t^\star$ be the true next token. Show that the derivative of the loss $-\log\Pr(\text{next token is }t^\star)$
with respect to $\logit(t)$ is $\Pr(\text{next token is }t)-1$ for $t=t^\star$ and $\Pr(\text{next token is }t)$
for every other $t$. What does training therefore do to the logits?
\end{exercise}

\begin{exercise}[Plausible or faithful]\label{ex:plausible}
A model is asked whether $3599$ is a prime number. It answers ``no''. When it is asked for its reason, it writes
``because $3599=59\times61$''. (a) Is this explanation plausible? (b) Describe an experiment that uses only inputs
and outputs and that could show that the explanation is not faithful. (c) Why can no experiment of this kind show
that the explanation is faithful?
\end{exercise}

\begin{exercise}[Three kinds of evidence]\label{ex:evidence}
Consider three statements about one neuron of a language model. (a) It is active on $93\%$ of all French tokens
and on $2\%$ of all other tokens. (b) When it is set to zero, the loss on French text rises and the loss on English
text stays the same. (c) Its outgoing weights directly raise the logits of common French words. Classify each
statement as observational, interventional or mechanistic. Then consider the claim ``the model uses this neuron to
recognise French''. For each statement, say which part of this claim it does not establish.
\end{exercise}

\begin{exercise}[A behaviour before a mechanism]\label{ex:behaviour}
In \cref{tab:predictions}, GPT-2 small gives \tok{Paris} a probability of $0.07$. Suppose that somebody presents
a detailed account of ``how GPT-2 small recalls that the Eiffel Tower is in Paris''. Which questions should the
audience ask first?
\end{exercise}

\begin{readinglist}
\item \citet{jacovi2020towards}, \emph{ACL}. The paper that introduced the distinction between plausible and
faithful explanations.
\item \citet{turpin2023language}, \emph{NeurIPS}. The experiments on explanations that leave out the true cause of
an answer.
\item \citet{bereska2024mechanistic}, \emph{TMLR}. A review of the concepts and methods of mechanistic
interpretability. It gives an overview of the topics of all eight sessions.
\end{readinglist}
